\documentclass[12pt,a4paper]{article}
\usepackage{../../../myconfig}
% Giữ nguyên đường dẫn cấu hình theo file mẫu của bạn.
% Đáp án và kiểm tra chỉ có trong comment, không in ra PDF.
\begin{document}
\titlebox{Chủ đề: Oxyz}{Phương trình đường thẳng}{Oxyz: Phương trình đường thẳng}
\begin{center}\textbf{Bài tập điểm danh -- 20 câu}\end{center}
\noindent Trong toàn bộ đề, các điểm và đường thẳng được xét trong không gian $Oxyz$. Mỗi câu chọn đúng một đáp án.\par\medskip
% Mức độ 1: Nhận biết

% Câu 1 | Dạng: Nhận biết vectơ chỉ phương | Mức độ: 1
% Nguồn: Chuyên đề 31, câu hỏi 5–6; câu 1. Biên tập cách hỏi/lựa chọn.
% Đáp án: C
% Kiểm tra: Vectơ chỉ phương là (2;-5;3).
\questionLinesManual{0}{1}{%
Cho đường thẳng $d:\begin{cases}x=3+2t\\y=4-5t\\z=-1+3t\end{cases}\ (t\in\mathbb R)$. Vectơ nào sau đây là một vectơ chỉ phương của $d$?
}{%
\begin{tasks}(4)
  \task $(2;-5;4)$
  \task $(3;-5;3)$
  \task $(2;-5;3)$
  \task $(2;-6;3)$
\end{tasks}
}

% Câu 2 | Dạng: Qua hai điểm | Mức độ: 1
% Nguồn: Chuyên đề 31, câu hỏi 5–6; câu 22. Biên tập cách hỏi/lựa chọn.
% Đáp án: A
% Kiểm tra: AB = (2;2;-2).
\questionLinesManual{0}{2}{%
Viết phương trình đường thẳng qua $A(1;0;1)$ và $B(3;2;-1)$.
}{%
\begin{tasks}(2)
  \task $\begin{cases}x=1+t\\y=t\\z=1-t\end{cases}\ (t\in\mathbb R)$
  \task $\begin{cases}x=1+t\\y=0\\z=1-t\end{cases}\ (t\in\mathbb R)$
  \task $\begin{cases}x=1+2t\\y=t\\z=2-t\end{cases}\ (t\in\mathbb R)$
  \task $\begin{cases}x=2+t\\y=t\\z=1-t\end{cases}\ (t\in\mathbb R)$
\end{tasks}
}

% Câu 3 | Dạng: Điểm và vectơ chỉ phương | Mức độ: 1
% Nguồn: Chuyên đề 31, câu hỏi 5–6; câu 25. Biên tập cách hỏi/lựa chọn.
% Đáp án: C
% Kiểm tra: Lấy điểm (2;0;-1) và vectơ (2;-3;1).
\questionLinesManual{0}{3}{%
Viết phương trình tham số của đường thẳng qua $M(2;0;-1)$, có vectơ chỉ phương $(2;-3;1)$.
}{%
\begin{tasks}(2)
  \task $\begin{cases}x=3+2t\\y=-3t\\z=-1+t\end{cases}\ (t\in\mathbb R)$
  \task $\begin{cases}x=2+3t\\y=-3t\\z=t\end{cases}\ (t\in\mathbb R)$
  \task $\begin{cases}x=2+2t\\y=-3t\\z=-1+t\end{cases}\ (t\in\mathbb R)$
  \task $\begin{cases}x=2+2t\\y=-4t\\z=-1+t\end{cases}\ (t\in\mathbb R)$
\end{tasks}
}

% Câu 4 | Dạng: Điểm thuộc đường thẳng | Mức độ: 1
% Nguồn: Chuyên đề 31, câu hỏi 5–6; câu 69. Biên tập cách hỏi/lựa chọn.
% Đáp án: C
% Kiểm tra: Thay t=2 được (1;-4;-5).
\questionLinesManual{0}{4}{%
Cho $d:\begin{cases}x=-1+t\\y=2-3t\\z=1-3t\end{cases}\ (t\in\mathbb R)$. Điểm nào sau đây thuộc $d$?
}{%
\begin{tasks}(4)
  \task $(2;-4;-5)$
  \task $(1;-5;-5)$
  \task $(1;-4;-5)$
  \task $(1;-4;-4)$
\end{tasks}
}

% Câu 5 | Dạng: Đổi dạng phương trình | Mức độ: 1
% Nguồn: Chuyên đề 31, câu hỏi 5–6; câu 23. Biên tập cách hỏi/lựa chọn.
% Đáp án: C
% Kiểm tra: \begin{cases}x=1+2t\\y=3t\\z=-2+t\end{cases}\ (t\in\mathbb R)
\questionLinesManual{0}{5}{%
Phương trình tham số của đường thẳng $d:\dfrac{x-1}{2}=\dfrac{y}{3}=\dfrac{z+2}{1}$ là
}{%
\begin{tasks}(2)
  \task $\begin{cases}x=1+3t\\y=3t\\z=-1+t\end{cases}\ (t\in\mathbb R)$
  \task $\begin{cases}x=2+2t\\y=3t\\z=-2+t\end{cases}\ (t\in\mathbb R)$
  \task $\begin{cases}x=1+2t\\y=3t\\z=-2+t\end{cases}\ (t\in\mathbb R)$
  \task $\begin{cases}x=1+2t\\y=2t\\z=-2+t\end{cases}\ (t\in\mathbb R)$
\end{tasks}
}
% Mức độ 2: Thông hiểu

% Câu 6 | Dạng: Tham số trong vectơ chỉ phương | Mức độ: 2
% Nguồn: Chuyên đề 31, câu hỏi 5–6; câu 17. Biên tập cách hỏi/lựa chọn.
% Đáp án: A
% Kiểm tra: u=2(2;1;2), nên a=b=4.
\questionLinesManual{0}{6}{%
Cho $d:\begin{cases}x=1+2t\\y=2+t\\z=-1+2t\end{cases}\ (t\in\mathbb R)$. Vectơ $\vec u=(a;2;b)$ là vectơ chỉ phương của $d$. Giá trị $a+b$ bằng
}{%
\begin{tasks}(4)
  \task $8$
  \task $10$
  \task $9$
  \task $7$
\end{tasks}
}

% Câu 7 | Dạng: Qua điểm, song song đường thẳng | Mức độ: 2
% Nguồn: Chuyên đề 31, câu hỏi 5–6; câu 66. Biên tập cách hỏi/lựa chọn.
% Đáp án: C
% Kiểm tra: Lấy hướng của d. Nếu A thuộc d, đường tìm được trùng d.
\questionLinesManual{0}{7}{%
Viết phương trình đường thẳng qua $A(2;1;-1)$ và có vectơ chỉ phương cùng phương với vectơ chỉ phương của $d:\begin{cases}x=1+t\\y=-1-2t\\z=2+t\end{cases}\ (t\in\mathbb R)$.
}{%
\begin{tasks}(2)
  \task $\begin{cases}x=3+t\\y=1-2t\\z=-1+t\end{cases}\ (t\in\mathbb R)$
  \task $\begin{cases}x=2+2t\\y=1-2t\\z=t\end{cases}\ (t\in\mathbb R)$
  \task $\begin{cases}x=2+t\\y=1-2t\\z=-1+t\end{cases}\ (t\in\mathbb R)$
  \task $\begin{cases}x=2+t\\y=1-3t\\z=-1+t\end{cases}\ (t\in\mathbb R)$
\end{tasks}
}

% Câu 8 | Dạng: Qua điểm, vuông góc mặt phẳng | Mức độ: 2
% Nguồn: Chuyên đề 31, câu hỏi 5–6; câu 40. Biên tập cách hỏi/lựa chọn.
% Đáp án: A
% Kiểm tra: Vectơ chỉ phương bằng pháp tuyến (2;-1;3).
\questionLinesManual{0}{8}{%
Viết phương trình đường thẳng qua $A(1;-2;3)$ và vuông góc với $(P):2x-y+3z+1=0$.
}{%
\begin{tasks}(2)
  \task $\begin{cases}x=1+2t\\y=-2-t\\z=3+3t\end{cases}\ (t\in\mathbb R)$
  \task $\begin{cases}x=1+2t\\y=-2-2t\\z=3+3t\end{cases}\ (t\in\mathbb R)$
  \task $\begin{cases}x=2+2t\\y=-2-t\\z=3+3t\end{cases}\ (t\in\mathbb R)$
  \task $\begin{cases}x=1+3t\\y=-2-t\\z=4+3t\end{cases}\ (t\in\mathbb R)$
\end{tasks}
}

% Câu 9 | Dạng: Trung tuyến tam giác | Mức độ: 2
% Nguồn: Ngân hàng PTT 2 (Phạm Thị Thủy); câu 262. Biên tập cách hỏi/lựa chọn.
% Đáp án: C
% Kiểm tra: Trung điểm BC là (2;-4;-4).
\questionLinesManual{0}{9}{%
Cho tam giác $ABC$ với $A(1;-3;4)$, $B(-2;-5;-7)$, $C(6;-3;-1)$. Viết phương trình trung tuyến từ $A$.
}{%
\begin{tasks}(2)
  \task $\begin{cases}x=1+2t\\y=-3-t\\z=5-8t\end{cases}\ (t\in\mathbb R)$
  \task $\begin{cases}x=2+t\\y=-3-t\\z=4-8t\end{cases}\ (t\in\mathbb R)$
  \task $\begin{cases}x=1+t\\y=-3-t\\z=4-8t\end{cases}\ (t\in\mathbb R)$
  \task $\begin{cases}x=1+t\\y=-3-2t\\z=4-8t\end{cases}\ (t\in\mathbb R)$
\end{tasks}
}

% Câu 10 | Dạng: Giao tuyến hai mặt phẳng | Mức độ: 2
% Nguồn: Ngân hàng PTT 2 (Phạm Thị Thủy); câu 257. Biên tập cách hỏi/lựa chọn.
% Đáp án: A
% Kiểm tra: Tích có hướng = (7;16;3), điểm trên giao tuyến (0;-\dfrac{1}{7};\dfrac{2}{7}).
\questionLinesManual{0}{10}{%
Cho $(P):x-y+3z-1=0$ và $(Q):3x-7z+2=0$. Vectơ nào là vectơ chỉ phương của giao tuyến?
}{%
\begin{tasks}(4)
  \task $(7;16;3)$
  \task $(7;15;3)$
  \task $(8;16;3)$
  \task $(7;16;4)$
\end{tasks}
}

% Câu 11 | Dạng: Giao điểm đường thẳng và mặt phẳng | Mức độ: 2
% Nguồn: Chuyên đề 31, câu hỏi 5–6; câu 87. Biên tập cách hỏi/lựa chọn.
% Đáp án: D
% Kiểm tra: Tham số trong phương trình đề bài: t=-1; giao điểm (1;1;0).
\questionLinesManual{0}{11}{%
Cho $d:\begin{cases}x=2+t\\y=-t\\z=3+3t\end{cases}\ (t\in\mathbb R)$ và $(P):x+y-z-2=0$. Tọa độ giao điểm là
}{%
\begin{tasks}(4)
  \task $(2;1;0)$
  \task $(1;0;0)$
  \task $(1;1;1)$
  \task $(1;1;0)$
\end{tasks}
}

% Câu 12 | Dạng: Chiếu điểm lên mặt phẳng | Mức độ: 2
% Nguồn: Chuyên đề 31, câu hỏi 5–6; câu 90. Biên tập cách hỏi/lựa chọn.
% Đáp án: A
% Kiểm tra: H=A-((n.A+d)/(n.n))n=(-4;-1;2).
\questionLinesManual{0}{12}{%
Hình chiếu vuông góc của $A(-4;5;2)$ lên $(P):y+1=0$ là
}{%
\begin{tasks}(4)
  \task $(-4;-1;2)$
  \task $(-3;-1;2)$
  \task $(-4;-2;2)$
  \task $(-4;-1;3)$
\end{tasks}
}

% Câu 13 | Dạng: Vị trí tương đối hai đường thẳng | Mức độ: 2
% Nguồn: Ngân hàng PTT 2 (Phạm Thị Thủy); câu 351. Biên tập cách hỏi/lựa chọn.
% Đáp án: D
% Kiểm tra: Hai vectơ cùng hướng; (3;1;-3) thuộc d1 (t=1), nên trùng nhau.
\questionLinesManual{0}{13}{%
Cho $d_1:\begin{cases}x=1+2t\\y=2-t\\z=-3\end{cases}\ (t\in\mathbb R)$ và $d_2:\begin{cases}x=3+2t\\y=1-t\\z=-3\end{cases}\ (t\in\mathbb R)$. Có bao nhiêu điểm chung của hai đường thẳng?
}{%
\begin{tasks}(4)
  \task Không có điểm chung
  \task Đúng một điểm chung
  \task Đúng hai điểm chung
  \task Vô số điểm chung
\end{tasks}
}
% Mức độ 3: Vận dụng

% Câu 14 | Dạng: Qua điểm, song song hai mặt phẳng | Mức độ: 3
% Nguồn: Chuyên đề 31, câu hỏi 5–6; câu 62. Biên tập cách hỏi/lựa chọn.
% Đáp án: A
% Kiểm tra: Lấy tích có hướng hai pháp tuyến: (1;0;-1).
\questionLinesManual{0}{14}{%
Viết phương trình đường thẳng qua $A(1;-2;3)$ có vectơ chỉ phương vuông góc với cả hai pháp tuyến của $(P):x+y+z+1=0$ và $(Q):x-y+z-2=0$.
}{%
\begin{tasks}(2)
  \task $\begin{cases}x=1+t\\y=-2\\z=3-t\end{cases}\ (t\in\mathbb R)$
  \task $\begin{cases}x=2+t\\y=-2\\z=3-t\end{cases}\ (t\in\mathbb R)$
  \task $\begin{cases}x=1+2t\\y=-2\\z=4-t\end{cases}\ (t\in\mathbb R)$
  \task $\begin{cases}x=1+t\\y=-2-t\\z=3-t\end{cases}\ (t\in\mathbb R)$
\end{tasks}
}

% Câu 15 | Dạng: Vuông góc hai hướng | Mức độ: 3
% Nguồn: Ngân hàng PTT 2 (Phạm Thị Thủy); câu 273. Biên tập cách hỏi/lựa chọn.
% Đáp án: B
% Kiểm tra: Lấy tích có hướng được (3;-4;-1).
\questionLinesManual{0}{15}{%
Viết phương trình đường thẳng qua $M(1;2;1)$, có vectơ chỉ phương vuông góc với cả $(1;1;-1)$ và $(2;1;2)$.
}{%
\begin{tasks}(2)
  \task $\begin{cases}x=1+3t\\y=2-5t\\z=1-t\end{cases}\ (t\in\mathbb R)$
  \task $\begin{cases}x=1+3t\\y=2-4t\\z=1-t\end{cases}\ (t\in\mathbb R)$
  \task $\begin{cases}x=1+4t\\y=2-4t\\z=2-t\end{cases}\ (t\in\mathbb R)$
  \task $\begin{cases}x=2+3t\\y=2-4t\\z=1-t\end{cases}\ (t\in\mathbb R)$
\end{tasks}
}

% Câu 16 | Dạng: Chiếu điểm lên đường thẳng | Mức độ: 3
% Nguồn: Ngân hàng PTT 2 (Phạm Thị Thủy); câu 278. Biên tập cách hỏi/lựa chọn.
% Đáp án: B
% Kiểm tra: Chân đường vuông góc H=(\dfrac{7}{3};-\dfrac{1}{3};-\dfrac{2}{3}).
\questionLinesManual{0}{16}{%
Cho $A(2;1;0)$ và $d:\begin{cases}x=1+2t\\y=-1+t\\z=-t\end{cases}\ (t\in\mathbb R)$. Hình chiếu vuông góc của A lên d là
}{%
\begin{tasks}(4)
  \task $(\dfrac{7}{3};-\dfrac{4}{3};-\dfrac{2}{3})$
  \task $(\dfrac{7}{3};-\dfrac{1}{3};-\dfrac{2}{3})$
  \task $(\dfrac{10}{3};-\dfrac{1}{3};-\dfrac{2}{3})$
  \task $(\dfrac{7}{3};-\dfrac{1}{3};\dfrac{1}{3})$
\end{tasks}
}

% Câu 17 | Dạng: Qua điểm, cắt và vuông góc đường thẳng | Mức độ: 3
% Nguồn: Chuyên đề 31, câu hỏi 7–8; câu 8. Biên tập cách hỏi/lựa chọn.
% Đáp án: D
% Kiểm tra: Chân đường vuông góc H=(2;1;1); lấy vectơ AH.
\questionLinesManual{0}{17}{%
Viết phương trình đường thẳng qua $A(1;0;2)$, cắt và vuông góc với $d:\begin{cases}x=1+t\\y=t\\z=-1+2t\end{cases}\ (t\in\mathbb R)$.
}{%
\begin{tasks}(2)
  \task $\begin{cases}x=1+2t\\y=t\\z=3-t\end{cases}\ (t\in\mathbb R)$
  \task $\begin{cases}x=2+t\\y=t\\z=2-t\end{cases}\ (t\in\mathbb R)$
  \task $\begin{cases}x=1+t\\y=0\\z=2-t\end{cases}\ (t\in\mathbb R)$
  \task $\begin{cases}x=1+t\\y=t\\z=2-t\end{cases}\ (t\in\mathbb R)$
\end{tasks}
}

% Câu 18 | Dạng: Khoảng cách điểm đến đường thẳng | Mức độ: 3
% Nguồn: Chuyên đề 31, câu hỏi 5–6; câu 104. Biên tập cách hỏi/lựa chọn.
% Đáp án: C
% Kiểm tra: d^2=56; chân H=(0;2;3).
\questionLinesManual{0}{18}{%
Khoảng cách từ $A(2;-4;-1)$ đến $d:\begin{cases}x=t\\y=2-t\\z=3+2t\end{cases}\ (t\in\mathbb R)$ bằng
}{%
\begin{tasks}(4)
  \task $\sqrt{57}$
  \task $\dfrac{1}{2}\sqrt{222}$
  \task $2\sqrt{14}$
  \task $\dfrac{1}{2}\sqrt{226}$
\end{tasks}
}
% Mức độ 4: Vận dụng kết hợp

% Câu 19 | Dạng: Qua điểm, cắt hai đường thẳng | Mức độ: 4
% Nguồn: Ngân hàng PTT 2 (Phạm Thị Thủy); câu 352. Biên tập cách hỏi/lựa chọn.
% Đáp án: B
% Kiểm tra: Giao điểm lần lượt (\dfrac{9}{2};-\dfrac{11}{2};10) và (-9;8;-14).
\questionLinesManual{0}{19}{%
Viết phương trình đường thẳng qua $M(0;-1;2)$ và cắt $d_1:\begin{cases}x=1+t\\y=-2-t\\z=3+2t\end{cases}\ (t\in\mathbb R)$, $d_2:\begin{cases}x=-1+2t\\y=4-t\\z=2+4t\end{cases}\ (t\in\mathbb R)$ tại hai điểm phân biệt.
}{%
\begin{tasks}(2)
  \task $\begin{cases}x=1+9t\\y=-1-9t\\z=2+16t\end{cases}\ (t\in\mathbb R)$
  \task $\begin{cases}x=9t\\y=-1-9t\\z=2+16t\end{cases}\ (t\in\mathbb R)$
  \task $\begin{cases}x=9t\\y=-1-10t\\z=2+16t\end{cases}\ (t\in\mathbb R)$
  \task $\begin{cases}x=10t\\y=-1-9t\\z=3+16t\end{cases}\ (t\in\mathbb R)$
\end{tasks}
}

% Câu 20 | Dạng: Đường vuông góc chung | Mức độ: 4
% Nguồn: Chuyên đề 31, câu hỏi 7–8; câu 18. Biên tập cách hỏi/lựa chọn.
% Đáp án: D
% Kiểm tra: Hai chân (1;0;0) và (2;1;-2); vectơ nối chân vuông góc cả hai hướng.
\questionLinesManual{0}{20}{%
Phương trình đường vuông góc chung của $d_1:\begin{cases}x=2+t\\y=1+t\\z=1+t\end{cases}\ (t\in\mathbb R)$ và $d_2:\begin{cases}x=t\\y=7-3t\\z=-t\end{cases}\ (t\in\mathbb R)$ là
}{%
\begin{tasks}(2)
  \task $\begin{cases}x=1+t\\y=0\\z=-2t\end{cases}\ (t\in\mathbb R)$
  \task $\begin{cases}x=2+t\\y=t\\z=-2t\end{cases}\ (t\in\mathbb R)$
  \task $\begin{cases}x=1+2t\\y=t\\z=1-2t\end{cases}\ (t\in\mathbb R)$
  \task $\begin{cases}x=1+t\\y=t\\z=-2t\end{cases}\ (t\in\mathbb R)$
\end{tasks}
}
\end{document}
